% !PRINTED 635
\weberpage{635}

In the discussion of this chapter we have, in accordance with the principal aim of the
investigation, regarded the coefficients $a_1, a_2, \ldots, a_n$ \gesperrt{as rational functions of}
$z$. Everything, however, remains entirely unchanged if we fix upon any \gesperrt{domain of
rationality} whatever for these quantities.

\begin{center}
\S\ 175. \gesperrt{Integral Functions of $z$.}
\end{center}

Every function $\omega$ of the field $\Omega$ satisfies, as we have seen, an equation of
lowest degree:

\weq{1}{\omega^e + b_1\omega^{e-1} + b_2\omega^{e-2} + \cdots + b_e = 0,}

whose coefficients $b_1, b_2, \ldots, b_e$ are rational functions of $z$.

1.\ \gesperrt{If these coefficients are integral rational functions of $z$, then $\omega$ is
called an integral algebraic function, or briefly an integral function, of $z$.}

For the integral algebraic functions the same theorems hold as for the integers (Bd.\ II,
\S\,149).

If we set

\weq{2}{\varphi(t) = t^e + b_1 t^{e-1} + b_2 t^{e-2} + \cdots + b_e,}

then, as we have seen in \S\,172,

\weq{3}{N(t - \omega) = \varphi(t)^f \quad (ef = n)}

is an integral power of $\varphi(t)$. If we therefore arrange $N(t-\omega)$ according to
powers of $t$, the coefficients all become integral rational functions again, and hence in
particular:

2.\ \gesperrt{The norm and the trace of an integral function $\omega$ are integral rational
functions of $z$.}

3.\ \gesperrt{A rational function of $z$ is an integral algebraic function only when it is
an integral rational function.}

For if $\omega = -b$ is rational, then $e = 1$ and $\omega + b = 0$ is the equation of
lowest degree for $\omega$, and hence $\omega$ is integral only when $b$ is integral and
rational.

4.\ \gesperrt{Every function $\eta$ in $\Omega$ can be transformed into an integral
algebraic function by multiplication by a rational function of $z$ different from zero.}

For every function $\eta$ satisfies an equation of lowest degree:

\weq{4}{b_0\eta^e + b_1\eta^{e-1} + \cdots + b_{e-1}\eta + b_e = 0,}
%CONT>>
